% Section 15.7, from the summary table of lectures/L15/appendix/a_theory.md and from
% lectures/L15/README.md: the objectives, the evaluation questions and the next lesson.

\section{Sammanfattning}
\label{c15:sec:sammanfattning}

\begin{booktable}{@{}p{12.5em}L@{}}
  \thead{Begrepp} & \thead{Formel} \\
  \midrule
  Rektangulär form & $z = x + jy$ \\
  Absolutbelopp & $\abs{z} = \sqrt{x^2 + y^2}$ \\
  Fasvinkel & $\delta = \arctan(y/x)$ + korrigering \\
  Polär form & $z = \abs{z}\angle\delta$ \\
  Polär till rektangulär form & $x = \abs{z}\cos\delta$, \ $y = \abs{z}\sin\delta$ \\
  Addition & $(x_1 + x_2) + j(y_1 + y_2)$ \\
\end{booktable}

\needspace{6\baselineskip}
\subsection*{Det här ska du kunna}
Efter det här kapitlet ska du:
\begin{itemize}
  \item Förstå vad komplexa tal är och hur de kan representeras i det komplexa talplanet.
  \item Kunna genomföra enkla beräkningar med komplexa tal på rektangulär form.
  \item Kunna omvandla mellan rektangulär och polär form samt ange absolutbelopp och fasvinkel.
\end{itemize}

\testadigsjalv
\begin{enumerate}
  \item Omvandla $z = 3 + j4$ till polär form. Ange absolutbeloppet och fasvinkeln i radianer.
  \item Beräkna $(2 + j3)(1 - j)$ på rektangulär form.
\end{enumerate}

\needspace{6\baselineskip}
\subsection*{Nästa kapitel}
\Kapref{16} fortsätter med komplexa tal: Eulers formel och komplexa tal för vektorberäkningar.
